\begin{proof}

\begin{enumerate}[label={$\langle1\rangle$\arabic*.},leftmargin=*]
\item Use the minimal resolution and effective
integral divisor of P4R.6. Push this
divisor to the minimal ruled model
$T_0$. Its intersection with a ruling
fibre is two. The adjunction and
Riemann--Hurwitz calculation in P4R.3
therefore gives exactly
\[
 D_0=2H+\sum_j a_jF_j,
 \qquad a_j\in\mathbb Z_{>0},
\]
with $H$ a smooth section. The reverse
blowup coefficient is the integral
number $\mu-1$. At a conductor crossing
one has $\mu\ge2$, so its exceptional
curve remains in the support and
subdivides the edge. At a smooth
conductor point the operation either
adds a rational leaf or leaves the
reduced curve unchanged. Thus
$A=\operatorname{Supp}D$ is connected,
is $H$ with rational trees attached,
and $\pi_1(A)\to\pi_1(T)$ is an
isomorphism.

\item Put $C=\operatorname{Sing}(S)_{\rm red}$.
The map $\rho$ is proper and is an
isomorphism outside $A$, and P4R.6
identifies $A$ with its full inverse
image of $C$. The continuous map
$T\cup_A C\to S$ is therefore a
compact-to-Hausdorff bijection and a
homeomorphism. The same cofibration
and van Kampen argument as P4R.3 gives
$\pi_1(S)\simeq\pi_1(C)$. This retains
the exceptional components that map
to points, as well as the strict
conductor components that map to
curves. P4R.1 and the normalization
graph calculation in P4R.3 give
$e(C)\ge0$. Finally P4.13 gives
$e(S_{\rm reg})\ge0$, so
$e(S)=e(S_{\rm reg})+e(C)\ge0$.

\end{enumerate}
\end{proof}
